% problem_id: S001-algebra-theorem-openness-flatness
% source_id: S001 (Stacks Project)
% source_file: algebra.tex
% statement_locator: algebra.tex lines 33406--33416
% proof_locator: algebra.tex lines 33418--33506
% env: theorem
% id_note: label 跨文件重名 ⇒ id 用文件限定形式
% status: complete (statement + author proof verbatim + reason + locators)
%
\begin{theorem}
\label{theorem-openness-flatness}
Let $R$ be a ring. Let $R \to S$ be a ring map of finite
presentation. Let $M$ be a finitely presented $S$-module.
The set
$$
\{ \mathfrak q \in \Spec(S) \mid
M_{\mathfrak q}\text{ is flat over }R\}
$$
is open in $\Spec(S)$.
\end{theorem}

\begin{proof}
Let $\mathfrak q \in \Spec(S)$ be a prime.
Let $\mathfrak p \subset R$ be the inverse image of $\mathfrak q$ in $R$.
Note that $M_{\mathfrak q}$ is flat over $R$ if and only if
it is flat over $R_{\mathfrak p}$.
Let us assume that $M_{\mathfrak q}$ is flat over $R$.
We claim that there exists a $g \in S$, $g \not \in \mathfrak q$
such that $M_g$ is flat over $R$.

\medskip\noindent
We first reduce to the case where $R$ and $S$ are
of finite type over $\mathbf{Z}$.
Choose a directed set $\Lambda$ and
a system $(R_\lambda \to S_\lambda, M_\lambda)$
as in Lemma \ref{lemma-limit-module-finite-presentation}.
Set $\mathfrak p_\lambda$ equal to the inverse image of
$\mathfrak p$ in $R_\lambda$.
Set $\mathfrak q_\lambda$ equal to the inverse image of
$\mathfrak q$ in $S_\lambda$.
Then the system
$$
((R_\lambda)_{\mathfrak p_\lambda},
(S_\lambda)_{\mathfrak q_\lambda},
(M_\lambda)_{\mathfrak q_{\lambda}})
$$
is a system as in
Lemma \ref{lemma-limit-module-essentially-finite-presentation}.
Hence by Lemma \ref{lemma-colimit-eventually-flat}
we see that for some $\lambda$ the module
$M_\lambda$ is flat over $R_\lambda$ at the prime
$\mathfrak q_{\lambda}$. Suppose we
can prove our claim for the system
$(R_\lambda \to S_\lambda, M_\lambda, \mathfrak q_{\lambda})$.
In other words, suppose that we can find a $g \in S_\lambda$,
$g \not\in \mathfrak q_\lambda$ such that $(M_\lambda)_g$
is flat over $R_\lambda$. By Lemma \ref{lemma-limit-module-finite-presentation}
we have $M = M_\lambda \otimes_{R_\lambda} R$ and hence
also $M_g = (M_\lambda)_g \otimes_{R_\lambda} R$. Thus by
Lemma \ref{lemma-flat-base-change} we deduce the claim
for the system $(R \to S, M, \mathfrak q)$.

\medskip\noindent
At this point we may assume that $R$ and $S$ are of finite type
over $\mathbf{Z}$. We may write $S$ as a quotient of a
polynomial ring $R[x_1, \ldots, x_n]$. Of course, we may replace
$S$ by $R[x_1, \ldots, x_n]$ and assume that $S$ is a polynomial
ring over $R$. In particular we see that $R \to S$ is flat
and all fibres rings $S \otimes_R \kappa(\mathfrak p)$
have global dimension $n$.

\medskip\noindent
Choose a resolution $F_\bullet$ of $M$ over $S$ with each
$F_i$ finite free, see Lemma \ref{lemma-resolution-by-finite-free}.
Let $K_n = \Ker(F_{n-1} \to F_{n-2})$. Note that
$(K_n)_{\mathfrak q}$ is flat over $R$, since each $F_i$
is flat over $R$ and by assumption on $M$, see Lemma
\ref{lemma-flat-ses}. In addition, the sequence
$$
0 \to
K_n/\mathfrak p K_n \to
F_{n-1}/ \mathfrak p F_{n-1} \to
\ldots \to
F_0 / \mathfrak p F_0 \to
M/\mathfrak p M \to
0
$$
is exact upon localizing at $\mathfrak q$, because of vanishing
of $\text{Tor}_i^{R_\mathfrak p}(\kappa(\mathfrak p), M_{\mathfrak q})$.
Since the global dimension of $S_\mathfrak q/\mathfrak p S_{\mathfrak q}$
is $n$ we conclude that $K_n / \mathfrak p K_n$ localized
at $\mathfrak q$ is a finite free module over
$S_\mathfrak q/\mathfrak p S_{\mathfrak q}$. By
Lemma \ref{lemma-free-fibre-flat-free} $(K_n)_{\mathfrak q}$
is free over $S_{\mathfrak q}$. In particular, there exists a
$g \in S$, $g \not \in \mathfrak q$ such that $(K_n)_g$
is finite free over $S_g$.

\medskip\noindent
By Lemma \ref{lemma-exact-on-fibres-open}
there exists a further localization $S_g$ such that
the complex
$$
0 \to K_n \to F_{n-1} \to \ldots \to F_0
$$
is exact on {\it all fibres} of $R \to S$. By
Lemma \ref{lemma-complex-exact-mod}
this implies that the cokernel of $F_1 \to F_0$ is
flat. This proves the theorem in the Noetherian case.
\end{proof}
